\subsection{Actual autarkic programs and their reusable proofs}

\begin{hbdefinition}{def:p12-har-autarkic-regime}{The source computation mechanism and this concrete instantiation}
Section~4.3 leaves the program representing square root unspecified,
while requiring its universal adequacy proof; it also leaves the
arithmetic-syntax interpreter and its free-variable bookkeeping to be
constructed. Here those programs are built using the chapter's primitive
recursion mechanism in the explicitly historical arithmetic regime of
\cref{def:p12-har-prime-regime,def:p12-har-prime-arithmetic}.
Thus \(\mathsf{nat}:\mathsf{Set}\), its constructor computation,
generated induction into \(\mathsf{Prop}\), propositional equality,
and the already proved order/semiring library are the assumptions.
The available classical DN axiom is not used in the following programs
or proofs. No proof is eliminated to select data.

The source's abstract \(\mathsf{nat}:\mathsf{Type}\) recursor
schemas describe the same primitive-recursion mechanism at that different
declared sort; they are not being silently substituted for this Set-based
library. Nor are these programs terms of the unchanged beta-only HOL
calculus. Each recursion below has a fixed admitted target and explicitly
typed branches. General universe/inductive extension consistency remains
separate. These are concrete source-based reconstructions, not programs
printed in the chapter or claims of Coq execution
\cite[pp.~1220--1223]{BarendregtGeuvers2001}.
\end{hbdefinition}

\begin{hbdefinition}{def:p12-har-autarkic-eqb}{The actual equality-code and data selector}
Use O and one as the false and true \emph{data codes}. Put
\[
\begin{aligned}
 z&:=\mathsf{nat\_rec}\ \mathsf{one}
       (\lambda j:\mathsf{nat}.\lambda r:\mathsf{nat}.O),\\
 b&:=\lambda n:\mathsf{nat}.\lambda q:\mathsf{nat}\to\mathsf{nat}.
       \lambda m:\mathsf{nat}.
          \mathsf{nat\_rec}\ O
            (\lambda j:\mathsf{nat}.\lambda r:\mathsf{nat}.q j)\ m,\\
 \mathsf{eqb}&:=\mathsf{nat\_rec}\ z\ b,\\
 \mathsf{choose}&:=\lambda c:\mathsf{nat}.\lambda u:\mathsf{nat}.
  \lambda v:\mathsf{nat}.
    \mathsf{nat\_rec}\ v
       (\lambda j:\mathsf{nat}.\lambda r:\mathsf{nat}.u)\ c.
\end{aligned}
\]
The outer eqb recursion has target \(\mathsf{nat}\to\mathsf{nat}\);
the two inner recursions and choose have target nat. Consequently
\(\mathsf{eqb}:\mathsf{nat}\to\mathsf{nat}\to\mathsf{nat}\) and
\(\mathsf{choose}:\mathsf{nat}\to\mathsf{nat}\to\mathsf{nat}\to\mathsf{nat}\).
Their actual computation equations are
\[
\begin{aligned}
 \mathsf{eqb}\ O\ O&\longrightarrow^*\mathsf{one},&
 \mathsf{eqb}\ O\ (S m)&\longrightarrow^*O,\\
 \mathsf{eqb}\ (S n)\ O&\longrightarrow^*O,&
 \mathsf{eqb}\ (S n)\ (S m)&\longrightarrow^*\mathsf{eqb}\ n\ m,\\
 \mathsf{choose}\ O\ u\ v&\longrightarrow^*v,&
 \mathsf{choose}\ (S c)\ u\ v&\longrightarrow^*u.
\end{aligned}
\]
The arrows expand these finite definitions and perform beta/iota
contractions; the equations are not postulated boolean axioms.
\end{hbdefinition}

\begin{hblemma}{thm:p12-har-autarkic-eqb}{proved=annex:p12.har.autarkic-eqb}{The complete equality-code specification}
For all \(n,m:\mathsf{nat}\),
\[
\begin{aligned}
 &\mathsf{eqb}\ n\ m=O\ \vee\ \mathsf{eqb}\ n\ m=\mathsf{one},\\
 &\mathsf{eqb}\ n\ m=\mathsf{one}\ \longleftrightarrow\ n=m,\\
 &\mathsf{eqb}\ n\ m=O\ \longleftrightarrow\ \neg(n=m).
\end{aligned}
\]
This is a constructed decision procedure, not classical excluded middle
used to define a Set-valued program.
\end{hblemma}

\begin{hblemma}{thm:p12-har-autarkic-squares}{proved=annex:p12.har.autarkic-squares}{Square growth and the unique integer square-root bounds}
Write \(Q(r):=\mathsf{times}\ r\ r\). Then
\[
 Q(S r)=Q(r)+S(r+r),\qquad Q(r)<Q(S r).
\]
Squaring preserves weak and strict order on nat. If two values r,s
satisfy \(Q(r)\le n<Q(S r)\) and \(Q(s)\le n<Q(S s)\),
then \(r=s\). These assertions follow from the proved arithmetic
library; no square-root existence assumption is used.
\end{hblemma}

\begin{hbdefinition}{def:p12-har-autarkic-square-root}{The omitted square-root program F, explicitly constructed}
The actual fixed-target recursion is
\[
 F:=\mathsf{nat\_rec}\ O\
  (\lambda n:\mathsf{nat}.\lambda r:\mathsf{nat}.
     \mathsf{choose}\ (\mathsf{eqb}\ (S n)\ (Q(S r)))\ (S r)\ r).
\]
Its step has type \(\mathsf{nat}\to\mathsf{nat}\to\mathsf{nat}\),
so \(F:\mathsf{nat}\to\mathsf{nat}\). Operationally it increments
the current root exactly when the next input reaches the next square.
All branching inputs are computed natural-number data.
For a meta-level natural number k, retain the source's unary numeral
representation, written here
\(\ulcorner k\urcorner_{\mathsf{nat}}:=S^k O\);
this is not a quotation of arithmetic syntax into the term datatype below.
The nat and T subscripts on quotation brackets are editorial sort tags
exposing the two different representation types.
\end{hbdefinition}

\begin{hbtheorem}{thm:p12-har-autarkic-square-root}{proved=annex:p12.har.autarkic-square-root}{The universal F adequacy proof and the actual 45 reduction}
The constructed program satisfies the source's precise proof obligation
\[
 \Pi n:\mathsf{nat}.\quad Q(F n)\le n<Q(S(F n)).
\]
Since \(S(F n)=F n+\mathsf{one}\), this is
\((F n)^2\le n<(F n+1)^2\). Its values are the unique natural-number
integer square roots specified by these bounds. In particular the
explicit definitions admit the finite computation
\[
 F\ulcorner45\urcorner_{\mathsf{nat}}
     \twoheadrightarrow_{\beta\iota\delta}
          \ulcorner6\urcorner_{\mathsf{nat}}.
\]
The annex supplies the invariant's branches and the finite computation
trace. No executed kernel replay is asserted by this mathematical
reduction derivation \cite[p.~1221]{BarendregtGeuvers2001}.
\end{hbtheorem}

\begin{hbdefinition}{def:p12-har-autarkic-syntax}{The actual arithmetic syntax and its five generated branches}
The source \(\mathrm{term}(\mathrm{Nat})\) contains terms over
\(\langle\mathrm{Nat};+,\times,0,1\rangle\); it is not a datatype
of arbitrary typed lambda terms. Make its omitted variable representation
explicit with the positive data declaration \(T:\mathsf{Set}\) and
constructors
\[
\begin{aligned}
 Z,U&:T,&\mathsf{Var}&:\mathsf{nat}\to T,\\
 \mathsf{Add},\mathsf{Mul}&:T\to T\to T.
\end{aligned}
\]
The numeral argument of Var is an index for a free variable name, not its
value. This is the historical Set-based instance of the source's stated
inductive syntax mechanism. For each fixed admitted target C, the generated
fold has seeds \(z,u:C\), variable branch \(v:\mathsf{nat}\to C\),
and two branches \(a,m:T\to T\to C\to C\to C\).
The parenthesized five-argument fold notation means five curried
applications, not a new tuple datatype. Its constructor computations are
\[
\begin{aligned}
 \mathsf{fold}_T^C(z,u,v,a,m)Z&\longrightarrow z,\\
 \mathsf{fold}_T^C(z,u,v,a,m)U&\longrightarrow u,\\
 \mathsf{fold}_T^C(z,u,v,a,m)(\mathsf{Var}\ i)&\longrightarrow v i,\\
 \mathsf{fold}_T^C(z,u,v,a,m)(\mathsf{Add}\ s\ t)&\longrightarrow
    a s t(\mathsf{fold}_T^C(z,u,v,a,m)s)(\mathsf{fold}_T^C(z,u,v,a,m)t),\\
 \mathsf{fold}_T^C(z,u,v,a,m)(\mathsf{Mul}\ s\ t)&\longrightarrow
    m s t(\mathsf{fold}_T^C(z,u,v,a,m)s)(\mathsf{fold}_T^C(z,u,v,a,m)t).
\end{aligned}
\]
Only direct recursive fields occur, so the source strictly positive
branch construction gives precisely the two recursive values/proofs.
For a fixed \(P:T\to\mathsf{Prop}\), the generated induction has
branches \(PZ\), \(PU\), \(\Pi i.P(\mathsf{Var}\ i)\), and
\(\Pi s,t.Ps\to Pt\to P(\mathsf{Add}\ s\ t)\), respectively
\(\Pi s,t.Ps\to Pt\to P(\mathsf{Mul}\ s\ t)\).
These are the particular source-licensed recursion/induction schemas,
not a parametricity assumption or a new top-sort-polymorphic constant
\cite[pp.~1206--1209,1221--1222]{BarendregtGeuvers2001}.
\end{hbdefinition}

\begin{hbdefinition}{def:p12-har-autarkic-interpreter}{The constructed valuation-sensitive interpreter and substitution}
For each fixed \(A:\mathsf{Set}\), take specified values
\(0_A,1_A:A\), operations \(+_A,\times_A:A\to A\to A\),
and a valuation \(\nu:\mathsf{nat}\to A\). Define the actual fold
\[
 E_\nu^A:=\mathsf{fold}_T^A
   (0_A,1_A,\nu,
    \lambda s,t:T.\lambda a,b:A.a+_A b,
    \lambda s,t:T.\lambda a,b:A.a\times_A b).
\]
Every branch has the type just declared; no algebraic laws are needed
to form this interpreter. Its original natural-number instance uses
\(A=\mathsf{nat}\), \(0_A=O\), \(1_A=\mathsf{one}\), and the
already constructed plus/times. Abbreviate that instance by \(E_\nu\).
The arbitrary fixed A version is an explicit interpretation schema,
not a claim that the chapter prints a polymorphic model interpreter.

Quotation of a finite arithmetic expression is defined by its five
clauses: 0 maps to Z, 1 to U, and a name \(x_i\) with meta-level
natural-number index i maps to
\(\mathsf{Var}\ulcorner i\urcorner_{\mathsf{nat}}\);
a sum maps to Add of the two quotations, and a product to Mul of them. Write it
as \(\ulcorner p\urcorner_T\). Under a valuation, \(p[\nu]\)
replaces every \(x_i\) by
\(\nu\ulcorner i\urcorner_{\mathsf{nat}}\) and uses the declared values
and operations. There are no binders in this arithmetic grammar.
Finite open contexts have an actual valuation extension as well. For a
fixed A, define the data selector
\[
 \mathsf{sel}_A:=\lambda c:\mathsf{nat}.\lambda u,v:A.
      \mathsf{nat\_rec}\ v
        (\lambda j:\mathsf{nat}.\lambda r:A.u)\ c.
\]
For distinct meta-level indices \(i_1,\ldots,i_k\) with given
well-typed values \(a_1,\ldots,a_k:A\) in a valid context, take
\[
 \nu_\Gamma(m):=
  \mathsf{sel}_A(\mathsf{eqb}\ m\ulcorner i_1\urcorner_{\mathsf{nat}})a_1
   \bigl(\cdots\mathsf{sel}_A
    (\mathsf{eqb}\ m\ulcorner i_k\urcorner_{\mathsf{nat}})a_k\,0_A\cdots\bigr).
\]
This finite nesting is a fully specified term (at meta-level k=0 it
is the constant \(0_A\) function); the annex proves each assigned-index
computation. It uses neither choice nor a proof-to-data map.
Also define actual syntactic substitution, for \(\sigma:\mathsf{nat}\to T\),
by
\[
 \mathsf{sub}_\sigma:=\mathsf{fold}_T^T
 (Z,U,\sigma,
    \lambda s,t:T.\lambda a,b:T.\mathsf{Add}\ a\ b,
    \lambda s,t:T.\lambda a,b:T.\mathsf{Mul}\ a\ b).
\]
Neither program consults a logical proof to return data.
\end{hbdefinition}

\begin{hbtheorem}{thm:p12-har-autarkic-interpreter}{proved=annex:p12.har.autarkic-interpreter}{Actual interpretation, quotation and open-valuations agreement}
For every finite expression in the declared arithmetic grammar and every
well-typed valuation in a valid context, expansion and constructor reduction
give
\[
 E_\nu^A\ulcorner p\urcorner_T
            \twoheadrightarrow_{\beta\iota\delta}p[\nu].
\]
Here \(p[\nu]\) is the A-valued interpreted expression, not a
purported raw Nat expression of type A. For the Nat instance this is
exactly the source arithmetic quotation/evaluation relationship, with
the missing free-variable valuation supplied. Moreover, if
\(\Pi i:\mathsf{nat}.\nu i=\mu i\) in primitive equality on A,
then \(\Pi t:T.E_\nu^A t=E_\mu^A t\). Finite-context agreement on
the indices actually occurring in a given quotation is sufficient too,
as the annex's constructor-by-constructor proof recipe demonstrates.
No function extensionality or parametricity theorem is assumed.
\end{hbtheorem}

\begin{hbtheorem}{thm:p12-har-autarkic-substitution}{proved=annex:p12.har.autarkic-substitution}{The full open-term substitution equation}
For every fixed A with the declared operations, every
\(\nu:\mathsf{nat}\to A\), \(\sigma:\mathsf{nat}\to T\), and
\(t:T\), put \(\nu_\sigma:=\lambda i:\mathsf{nat}.E_\nu^A(\sigma i)\).
Then
\[
 E_\nu^A(\mathsf{sub}_\sigma t)=E_{\nu_\sigma}^A t.
\]
The equality is proved by the actual five generated induction branches,
not by supposing interpretation automatically commutes with substitution.
\end{hbtheorem}

\begin{hblemma}{thm:p12-har-autarkic-prime-bound}{proved=annex:p12.har.autarkic-prime-bound}{The actual finite coefficient bound needed by primality search}
For \(\mathsf{one}\le d\), every coefficient satisfies
\(x\le xd\). For such a d, consequently,
\[
 \mathsf{div}\ d\ n\to
      \exists x:\mathsf{nat}.(x\le n)\wedge(xd=n).
\]
In particular the source's proper-divisor proposition has the finite
search characterization
\[
\begin{aligned}
 &\exists d:\mathsf{nat}.\mathsf{propdiv}\ d\ n\\
 &\quad\longleftrightarrow\quad
 \exists d,x:\mathsf{nat}.
  (d\le n)\wedge(x\le n)\wedge(\mathsf{one}<d)\wedge(d<n)\wedge(xd=n).
\end{aligned}
\]
Both existentials are in Prop. This bounds the later data-search loops
used by the actual KPrime program in \cref{def:p12-har-kprime-program,thm:p12-har-kprime-correct}.
That program and its characteristic proof are supplied separately;
the loops are not obtained by extracting these Prop witnesses into data.
\end{hblemma}

\begin{proofboundary}{def:p12-har-autarkic-residual}{What these constructions do not supply}
The source square-root program and its full reusable proof are now supplied,
as are the actual arithmetic-syntax interpreter and its open-term laws.
The actual KPrime term and all order/boolean/bounded-fold/characteristic
proofs are supplied in \cref{def:p12-har-kprime-program,thm:p12-har-kprime-correct}.
The actual polynomial program and universal interpretation proof are
supplied in \cref{def:p12-har-poly-normalizer,def:p12-har-poly-quotation,thm:p12-har-poly-preservation}.
Its nontrivial source-pattern output is given with the explicit
association convention in \cref{thm:p12-har-poly-example}; it is not
identity G or a claim of syntax equality across different parses.
The external primitive-recursive representation templates and the distinct
internal formula compiler are supplied in
\cref{thm:p12-har-prcompiler-representation,thm:p12-har-prcompiler-ok,thm:p12-har-prcompiler-predicate}.
They are not consequences of these particular programs or an automatic
inverse from arbitrary propositions to formula data. The Pocklington/Fermat
variant and large-number certificate remain distinct additional constructions.
The global
extension metatheory and arithmetic sequence-coding bridge remain separate.
\end{proofboundary}
